\chapter{Kinematic and Dynamic Modeling}\label{ch:kinematic-and-dynamic-modeling}

\section{Coordinate Frames and Generalized Coordinates}\label{coordinate-frames-and-generalized-coordinates}

Let \{I\} denote an inertial frame fixed to the level ground and \{B\} a
body frame whose origin lies on the common wheel-center line. The body x
axis is forward, y is along the wheel axis, and z is upward at the
upright equilibrium. The planar pose is \ensuremath{\eta}=[X,Y,\ensuremath{\psi}]\ensuremath{^{T}}, where \ensuremath{\psi} is yaw.
The body twist is \ensuremath{\nu}=[v\_x,v\_y,\ensuremath{\dot{\psi}}]\ensuremath{^{T}}. The pitch angle \ensuremath{\theta} describes
rotation of the body about the common wheel axis, with \ensuremath{\theta}=0 at the
selected upright reference. The generalized coordinates are
q=[X,Y,\ensuremath{\psi},\ensuremath{\theta}]\ensuremath{^{T}}.

\begin{equation}
{}^{I}\mathbf{v}=\mathbf{R}(\psi){}^{B}\mathbf{v},\qquad
\mathbf{R}(\psi)=\begin{bmatrix}\cos\psi&-\sin\psi\\\sin\psi&\cos\psi\end{bmatrix}
\label{eq:6-1}
\end{equation}

Equation (6.1) transforms the body-frame translational velocity to
inertial coordinates. The sign of \ensuremath{\psi} follows the right-hand rule about
z\_I. All controller and plot signals should use radians internally and
state degrees only when presenting results.

\begin{longtable}[]{@{}
  >{\raggedright\arraybackslash}p{(\columnwidth - 0\tabcolsep) * \real{1.0000}}@{}}
\toprule\noalign{}
\begin{minipage}[b]{\linewidth}\raggedright
\textless INSERT FIGURE 6.1: Draw the inertial and body frames, pose
X-Y-psi, body velocities vx and vy, pitch theta, center of mass,
gravity, and common wheel axis. Show the upright reference and positive
rotations\textgreater{}\\
What the figure should show: Draw the inertial and body frames, pose
X-Y-psi, body velocities vx and vy, pitch theta, center of mass,
gravity, and common wheel axis. Show the upright reference and positive
rotations\\
Likely source: MATLAB vector graphic based on final CAD
coordinates\strut
\end{minipage} \\
\midrule\noalign{}
\endhead
\bottomrule\noalign{}
\endlastfoot
\end{longtable}

\begin{center}\singlespacing
\phantomsection\addcontentsline{lof}{figure}{\protect\numberline{6.1}Inertial and body coordinate frames used by the model.}
\textbf{Figure 6.1. Inertial and body coordinate frames used by the model.}
\end{center}

\section{Collinear Mecanum-Wheel Kinematics}\label{collinear-mecanum-wheel-kinematics}

\subsection{Wheel contact velocity}\label{wheel-contact-velocity}

Wheel i is located at p\_i=[x\_i,y\_i]\ensuremath{^{T}} in \{B\}. The planar
velocity of its center is the body translation plus the yaw
contribution.

\begin{equation}
\mathbf{v}_{i}^{B}=\begin{bmatrix}v_x\\v_y\end{bmatrix}+
\dot{\psi}\begin{bmatrix}-y_i\\x_i\end{bmatrix}
\label{eq:6-2}
\end{equation}

For the ideal collinear arrangement x\_i=0. The Mecanum roller
constraint is represented by a sign s\_i\ensuremath{\in}\{\ensuremath{-}1,+1\} for an ideal
45-degree roller. The sign is defined operationally: it specifies
whether positive wheel rotation couples positively or negatively to
v\_y. This avoids relying on wheel labels such as left-handed or
right-handed, whose interpretation can change with viewing direction.

\subsection{Inverse kinematics}\label{inverse-kinematics}

With effective wheel radius r and the stated sign convention, the ideal
inverse kinematics reduce to Equation (6.3). Each row states the wheel
rate needed for the local longitudinal velocity v\_x\ensuremath{-}y\_i\ensuremath{\dot{\psi}} and the
lateral component selected by s\_i.

\begin{equation}
\boldsymbol{\omega}=\frac{1}{r}\mathbf{H}\boldsymbol{\nu},\qquad
\mathbf{H}=\begin{bmatrix}1&s_1&-y_1\\1&s_2&-y_2\\1&s_3&-y_3\\1&s_4&-y_4\end{bmatrix},\quad
\boldsymbol{\nu}=\begin{bmatrix}v_x&v_y&\dot{\psi}\end{bmatrix}^{T}
\label{eq:6-3}
\end{equation}

The matrix H has dimensions 4\ensuremath{\times}3. The first column corresponds to
longitudinal velocity, the second to lateral velocity, and the third to
yaw. The equation is symbolic until r, y\_i, and s\_i are measured. If
the roller angle differs materially from 45 degrees, the second column
must be scaled from the actual constraint geometry rather than retaining
\ensuremath{\pm}1.

\begin{longtable}[]{@{}
  >{\raggedright\arraybackslash}p{(\columnwidth - 0\tabcolsep) * \real{1.0000}}@{}}
\toprule\noalign{}
\begin{minipage}[b]{\linewidth}\raggedright
\textless INSERT FIGURE 6.2: For each wheel show positive spin, roller
axis, effective constrained force direction, sign si, axial position yi,
and the local velocity contribution from positive yaw. Use the installed
wheel order\textgreater{}\\
What the figure should show: For each wheel show positive spin, roller
axis, effective constrained force direction, sign si, axial position yi,
and the local velocity contribution from positive yaw. Use the installed
wheel order\\
Likely source: MATLAB vector diagram checked against the physical
wheels\strut
\end{minipage} \\
\midrule\noalign{}
\endhead
\bottomrule\noalign{}
\endlastfoot
\end{longtable}

\begin{center}\singlespacing
\phantomsection\addcontentsline{lof}{figure}{\protect\numberline{6.2}Idealized force direction for each collinear Mecanum wheel.}
\textbf{Figure 6.2. Idealized force direction for each collinear Mecanum wheel.}
\end{center}

\subsection{Forward kinematics and rank}\label{forward-kinematics-and-rank}

When H has full column rank, the least-squares body twist follows from
the Moore-Penrose pseudoinverse.

\begin{equation}
\boldsymbol{\nu}=r\left(\mathbf{H}^{T}\mathbf{H}\right)^{-1}\mathbf{H}^{T}\boldsymbol{\omega}
=r\mathbf{H}^{\dagger}\boldsymbol{\omega},\qquad \operatorname{rank}(\mathbf{H})=3
\label{eq:6-4}
\end{equation}

Full planar mobility requires rank(H)=3. Both roller signs must be
present; otherwise the longitudinal and lateral columns are dependent.
In addition, the axial-position column must not be expressible as a
linear combination of the constant and sign columns. For four wheels, a
practical condition is that at least one roller-sign group contains
wheels at distinct y positions while both signs are represented. These
conditions must be evaluated for the measured order rather than assumed.

\par\medskip\noindent\singlespacing
\phantomsection\addcontentsline{lot}{table}{\protect\numberline{6.1}Mobility-rank conditions for the collinear wheel arrangement}
\textbf{Table 6.1. Mobility-rank conditions for the collinear wheel arrangement}\par

\begin{longtable}[]{@{}
  >{\raggedright\arraybackslash}p{(\columnwidth - 4\tabcolsep) * \real{0.4375}}
  >{\raggedright\arraybackslash}p{(\columnwidth - 4\tabcolsep) * \real{0.1562}}
  >{\raggedright\arraybackslash}p{(\columnwidth - 4\tabcolsep) * \real{0.4062}}@{}}
\toprule\noalign{}
\begin{minipage}[b]{\linewidth}\raggedright
\textbf{Geometry or roller case}
\end{minipage} & \begin{minipage}[b]{\linewidth}\raggedright
\textbf{Expected rank}
\end{minipage} & \begin{minipage}[b]{\linewidth}\raggedright
\textbf{Consequence}
\end{minipage} \\
\midrule\noalign{}
\endhead
\bottomrule\noalign{}
\endlastfoot
Both signs present and axial positions provide an independent yaw column
& 3 & Independent ideal vx, vy, and yaw commands are kinematically
available. \\
All wheels have the same roller sign & \ensuremath{\leq}2 & Longitudinal and lateral
columns are dependent; full planar mobility is lost. \\
Each sign occurs only at one repeated axial location & \ensuremath{\leq}2 & Yaw column
can become a combination of longitudinal and lateral columns. \\
Wheel locations nearly create dependence & 3 but ill-conditioned & Large
wheel speeds\slash{}torques and noise amplification; evaluate condition
number. \\
\end{longtable}

\section{Wheel Torque and Body Wrench}\label{wheel-torque-and-body-wrench}

Power consistency maps wheel torques to the body wrench
w=[F\_x,F\_y,M\_z]\ensuremath{^{T}}. If wheel-speed losses are neglected, \ensuremath{\tau}\ensuremath{^{T}}\ensuremath{\omega}=w\ensuremath{^{T}}\ensuremath{\nu},
giving Equation (6.5).

\begin{equation}
\mathbf{w}=\begin{bmatrix}F_x&F_y&M_z\end{bmatrix}^{T}
=\frac{1}{r}\mathbf{H}^{T}\boldsymbol{\tau}
\label{eq:6-5}
\end{equation}

Because four actuators produce three planar wrench components, one
internal null-space direction remains when H has rank three. A weighted
minimum-effort allocation is given by Equation (6.6), where W is
positive definite and can penalize actuators differently.

\begin{equation}
\boldsymbol{\tau}^{\star}=r\mathbf{W}^{-1}\mathbf{H}
\left(\mathbf{H}^{T}\mathbf{W}^{-1}\mathbf{H}\right)^{-1}\mathbf{w}_{d}
\label{eq:6-6}
\end{equation}

The allocator must retain torque headroom for pitch stabilization. A
planar command that uses all available current can prevent recovery from
a pitch disturbance. A practical implementation therefore prioritizes or
reserves the longitudinal collective torque needed by the balance loop,
then allocates lateral and yaw demands subject to individual current,
speed, voltage, and thermal limits. Null-space torque should normally be
zero unless it serves a documented objective, because internal wheel
forces can increase slip and power loss without changing the ideal body
wrench.

\section{Nonlinear Dynamic Model}\label{nonlinear-dynamic-model}

The general model uses q=[X,Y,\ensuremath{\psi},\ensuremath{\theta}]\ensuremath{^{T}} and four wheel torques. The
Euler-Lagrange form is shown in Equation (6.7). M(q) is the mass matrix;
C(q,\ensuremath{\dot{q}}) contains velocity-dependent terms; g(q) is the gravity vector; D
represents dissipative effects; B(q) maps actuator torques; and d
contains unmodeled contact, slip, and disturbance forces.

\begin{equation}
\mathbf{M}(\mathbf{q})\ddot{\mathbf{q}}+
\mathbf{C}(\mathbf{q},\dot{\mathbf{q}})\dot{\mathbf{q}}+
\mathbf{g}(\mathbf{q})+\mathbf{D}\dot{\mathbf{q}}
=\mathbf{B}(\mathbf{q})\boldsymbol{\tau}+\mathbf{d}
\label{eq:6-7}
\end{equation}

A complete numerical M matrix requires the final component masses,
center-of-mass coordinates, and inertia tensors. Until those data exist,
the model is kept symbolic. The principal assumptions are a rigid body,
a rigid level ground plane, continuous wheel contact, one dominant
unstable pitch coordinate, known wheel geometry, and an idealized roller
constraint. The initial model neglects detailed roller inertia,
intermittent roller impacts, tire deformation, bearing friction
nonlinearities, battery voltage sag, sensor noise, and CAN jitter. Each
omitted effect is either included as a disturbance in sensitivity
studies or added after hardware data show that it materially affects the
result.

\par\medskip\noindent\singlespacing
\phantomsection\addcontentsline{lot}{table}{\protect\numberline{6.2}Kinematic and dynamic symbols}
\textbf{Table 6.2. Kinematic and dynamic symbols}\par

\begin{longtable}[]{@{}
  >{\raggedright\arraybackslash}p{(\columnwidth - 4\tabcolsep) * \real{0.1538}}
  >{\raggedright\arraybackslash}p{(\columnwidth - 4\tabcolsep) * \real{0.6615}}
  >{\raggedright\arraybackslash}p{(\columnwidth - 4\tabcolsep) * \real{0.1846}}@{}}
\toprule\noalign{}
\begin{minipage}[b]{\linewidth}\raggedright
\textbf{Symbol}
\end{minipage} & \begin{minipage}[b]{\linewidth}\raggedright
\textbf{Definition}
\end{minipage} & \begin{minipage}[b]{\linewidth}\raggedright
\textbf{Unit}
\end{minipage} \\
\midrule\noalign{}
\endhead
\bottomrule\noalign{}
\endlastfoot
X, Y & Inertial position of body-frame origin & m \\
\ensuremath{\psi} & Yaw angle & rad \\
\ensuremath{\theta} & Pitch angle from upright & rad \\
v\_x, v\_y & Body-frame translation velocities & m\slash{}s \\
r & Effective loaded wheel radius & m \\
y\_i & Wheel-i position along common axis & m \\
s\_i & Wheel-i roller sign & --- \\
\ensuremath{\omega}\_i & Wheel-i angular speed & rad\slash{}s \\
\ensuremath{\tau}\_i & Wheel-i output torque applied in positive wheel direction &
N\ensuremath{\cdot}m \\
m\_b & Body mass participating in pitch motion & kg \\
m\_t & Equivalent translated mass & kg \\
\ensuremath{\ell} & Distance from wheel axis to body COM & m \\
I\_p & Body pitch inertia about COM & kg\ensuremath{\cdot}m\ensuremath{^{2}} \\
b\_x, b\_\ensuremath{\theta} & Equivalent longitudinal and pitch damping & N\ensuremath{\cdot}s\slash{}m,
N\ensuremath{\cdot}m\ensuremath{\cdot}s\slash{}rad \\
\end{longtable}

\section{Reduced Longitudinal Balance Model}\label{reduced-longitudinal-balance-model}

For initial controller design, longitudinal motion and pitch are
isolated from lateral and yaw dynamics. The body center of mass is
located a distance \ensuremath{\ell} from the wheel axis. The equivalent translated mass
m\_t includes the body and reflected wheel\slash{}rotor contributions used by
the chosen derivation. The collective wheel torque is \ensuremath{\tau}\_\ensuremath{\Sigma}=\ensuremath{\Sigma}\ensuremath{\tau}\_i for the
wheel combination that produces positive forward motion. With positive \ensuremath{\theta}
defined as forward lean and positive wheel torque producing forward
ground force, one consistent reduced model is Equation (6.8).

\begin{equation}
\begin{bmatrix}m_t&m_b\ell\cos\theta\\m_b\ell\cos\theta&I_p+m_b\ell^2\end{bmatrix}
\begin{bmatrix}\ddot{x}\\\ddot{\theta}\end{bmatrix}+
\begin{bmatrix}-m_b\ell\sin\theta\,\dot{\theta}^{2}\\-m_bg\ell\sin\theta\end{bmatrix}+
\begin{bmatrix}b_x\dot{x}\\b_\theta\dot{\theta}\end{bmatrix}
=\begin{bmatrix}1\slash{}r\\-1\end{bmatrix}\tau_{\Sigma}
\label{eq:6-8}
\end{equation}

The input vector contains the forward force \ensuremath{\tau}\_\ensuremath{\Sigma}/r and the opposing
reaction torque \ensuremath{-}\ensuremath{\tau}\_\ensuremath{\Sigma} on the body. The equation assumes that the four
wheel torques contributing to forward balance have already been
transformed into a consistent sign convention. If the actuator reports
motor-side rather than output torque, gear ratio and efficiency must be
applied explicitly. The model should be checked by energy and
dimensional consistency before numerical use.

\begin{longtable}[]{@{}
  >{\raggedright\arraybackslash}p{(\columnwidth - 0\tabcolsep) * \real{1.0000}}@{}}
\toprule\noalign{}
\begin{minipage}[b]{\linewidth}\raggedright
\textless INSERT FIGURE 6.3: Draw the reduced body and wheel system with
body mass mb, COM distance l, pitch inertia Ip, wheel radius r, gravity,
ground force, collective actuator torque, body reaction torque, x
displacement, and positive theta\textgreater{}\\
What the figure should show: Draw the reduced body and wheel system with
body mass mb, COM distance l, pitch inertia Ip, wheel radius r, gravity,
ground force, collective actuator torque, body reaction torque, x
displacement, and positive theta\\
Likely source: author-created free-body diagram verified against the
sign convention\strut
\end{minipage} \\
\midrule\noalign{}
\endhead
\bottomrule\noalign{}
\endlastfoot
\end{longtable}

\begin{center}\singlespacing
\phantomsection\addcontentsline{lof}{figure}{\protect\numberline{6.3}Reduced inverted-pendulum free-body representation.}
\textbf{Figure 6.3. Reduced inverted-pendulum free-body representation.}
\end{center}

\section{Linearization and State-Space Form}\label{linearization-and-state-space-form}

Near \ensuremath{\theta}=0, sin\ensuremath{\theta}\ensuremath{\approx}\ensuremath{\theta} and cos\ensuremath{\theta}\ensuremath{\approx}1. Let E be the 2\ensuremath{\times}2 mass matrix evaluated at
the upright equilibrium, K contain the linear gravitational term, and D
contain the selected viscous damping. Defining x=[x,\ensuremath{\theta},\ensuremath{\dot{x}},\ensuremath{\dot{\theta}}]\ensuremath{^{T}} gives
the continuous-time model in Equation (6.9).

\begin{equation}
\dot{\mathbf{x}}=\mathbf{A}\mathbf{x}+\mathbf{B}u,\qquad
\mathbf{x}=\begin{bmatrix}x&\theta&\dot{x}&\dot{\theta}\end{bmatrix}^{T},\qquad
\mathbf{A}=\begin{bmatrix}\mathbf{0}&\mathbf{I}\\
\mathbf{E}^{-1}\mathbf{K}&-\mathbf{E}^{-1}\mathbf{D}\end{bmatrix}
\label{eq:6-9}
\end{equation}

The exact A and B entries should be generated symbolically from the
adopted parameter definitions to avoid transcription errors. The model
is discretized at the measured controller interval, preferably using a
zero-order hold. Controllability must be checked with the final
numerical parameters using Equation (6.10); observability must likewise
be checked for the selected measured outputs.

\begin{equation}
\mathcal{C}=\begin{bmatrix}\mathbf{B}&\mathbf{A}\mathbf{B}&\mathbf{A}^{2}\mathbf{B}&\mathbf{A}^{3}\mathbf{B}\end{bmatrix},
\qquad \operatorname{rank}(\mathcal{C})=4\;\text{required}
\label{eq:6-10}
\end{equation}

\par\medskip\noindent\singlespacing
\phantomsection\addcontentsline{lot}{table}{\protect\numberline{6.3}Model parameters and identification method}
\textbf{Table 6.3. Model parameters and identification method}\par

\begin{longtable}[]{@{}
  >{\raggedright\arraybackslash}p{(\columnwidth - 4\tabcolsep) * \real{0.2388}}
  >{\raggedright\arraybackslash}p{(\columnwidth - 4\tabcolsep) * \real{0.3806}}
  >{\raggedright\arraybackslash}p{(\columnwidth - 4\tabcolsep) * \real{0.3806}}@{}}
\toprule\noalign{}
\begin{minipage}[b]{\linewidth}\raggedright
\textbf{Parameter}
\end{minipage} & \begin{minipage}[b]{\linewidth}\raggedright
\textbf{Source for final value}
\end{minipage} & \begin{minipage}[b]{\linewidth}\raggedright
\textbf{Uncertainty or validation method}
\end{minipage} \\
\midrule\noalign{}
\endhead
\bottomrule\noalign{}
\endlastfoot
m\_b and m\_t & Physical component weights and final CAD grouping &
Scale accuracy; compare CAD and assembled mass \\
\ensuremath{\ell} and COM offsets & CAD mass properties and physical balance\slash{}suspension
test & Repeat after wiring or battery changes \\
I\_p and I\_z & CAD plus pendulum\slash{}system-identification test & Report
method and confidence interval \\
r & Loaded rollout test & Repeat on test surface and at representative
load \\
y\_i and s\_i & Dimensioned CAD and wheel-direction inspection &
Cross-check against single-wheel command test \\
Actuator delay\slash{}time constant & Command\slash{}feedback step test & Fit model;
report sample rate and confidence \\
Torque\slash{}current limits & Installed actuator configuration and measured
current & Distinguish manufacturer rating from selected limit \\
b\_x and b\_\ensuremath{\theta} & Coast-down or fitted response & Treat as uncertain in
sensitivity analysis \\
Friction\slash{}slip parameters & Dedicated velocity\slash{}force experiments &
Surface-dependent; do not overfit one trial \\
\end{longtable}

\section{Model Limitations}\label{model-limitations}

The ideal model predicts achievable velocity and generalized wrench, not
exact ground forces at every roller. Mecanum wheels experience discrete
changes in the contacting roller, varying normal loads, and slip
\cite{bayar2020contact}--\cite{barcenas2025robust}. The body may also roll or flex even when the
principal unstable motion is pitch. The actuator\textquotesingle s
internal current and velocity loops, CAN command update, gearbox
backlash, and encoder location modify the commanded-to-wheel response.
These effects should be added only at a level supported by measurement.

The purpose of the reduced model is to make controller assumptions
visible and testable. Agreement should be evaluated over multiple
initial conditions and command directions. If one parameter set cannot
explain both small-signal balance and larger motion trials, the
residuals should guide model refinement rather than being hidden by
retuning the controller.
